\section{Stack tecnológica}
\subsection{Backend}
\begin{itemize}[nosep]
  \item Java 23, Spring Boot 3.3.x
  \item Spring Cloud Gateway (API Gateway)
  \item Spring Data JPA + PostgreSQL 16 (persistência)
  \item Flyway (migrações)
  \item RabbitMQ (AMQP) para eventos assíncronos
  \item SpringDoc OpenAPI/Swagger UI
  \item Resilience4j (Circuit Breaker/Retry) no \texttt{sync-service} para integrações externas
  \item Integrações: Stripe (pagamentos/webhooks), Cloudinary (media), Ably (chat/webhooks)
\end{itemize}

\subsection{Frontend}
\begin{itemize}[nosep]
  \item Next.js (App Router) + TypeScript
  \item Bun (runtime + package manager)
  \item Tailwind CSS
  \item Axios (camada de comunicação com a API via Gateway)
  \item Playwright (testes E2E/integração)
\end{itemize}

\subsection{Infraestrutura}
\begin{itemize}[nosep]
  \item Docker + Docker Compose (orquestração local/deploy)
  \item Postgres (container único com múltiplas bases de dados)
  \item RabbitMQ (broker + management)
\end{itemize}

\subsection{Justificação das Escolhas}
\begin{itemize}
    \item \textbf{Java (Spring Boot) vs. Node.js:} Optámos por Java no backend devido ao nosso maior nível de conhecimento e controlo sobre a linguagem para gerir a lógica complexa dos microserviços.
    \item \textbf{Next.js vs. Express/React nativo:} Embora tenhamos experiência com Express, escolhemos o Next.js pela enorme facilidade de uso, convenções integradas de roteamento e otimização imediata.
    \item \textbf{TypeScript vs. JavaScript:} Escolhemos TypeScript em todo o frontend porque a tipagem estática previne erros em \textit{runtime} e melhora a experiência de desenvolvimento e manutenção.
    \item \textbf{Tailwind CSS vs. CSS Tradicional:} O Tailwind foi selecionado porque a sua abordagem \textit{utility-first} acelera drasticamente a criação da UI sem termos de sair dos ficheiros dos componentes.
\end{itemize}